%
% claude-session-end-unordered.tex  (FROZEN NEGATIVE CONTROL: do not edit)
%
% NewHello is enabled while the old end is in progress, and ReleaseTty frees the name by name.
%
% Z Specification for how a Claude session ends when biffd owns its record:
% EndOwnSession (a connection of the current run ends its own session),
% ExpireClaude (retention removes a session nobody is attached to), the
% ordering of a clear (/clear ends the old session before the new hello), and the
% obligation that no row outlives an end plus retention.
%
% Companion to session-model.tex (SessionEnd, RetentionExpire for CLI sessions,
% the daemon's records and runs) and to sentinel-reap.tex (BadOrphan, findings G1
% and G2 of the sentinel mechanism).  It is kept apart because those two
% specifications already have large state spaces, and this one needs the effects
% of an end (the relay row, the tty-name reservation, the logout event, the held
% mail), which neither carries.
%
% Type-check: fuzz -t claude-session-end.tex
% Model-check: see the end of this file.
%

\documentclass[a4paper,10pt,fleqn]{article}
\usepackage[margin=1in]{geometry}
\usepackage{fuzz}
\usepackage[colorlinks=true,linkcolor=blue,citecolor=blue,urlcolor=blue]{hyperref}

\begin{document}

\title{Biff Claude Session End: A Z Specification}
\author{Punt Labs}
\date{November 2026}
\maketitle

%%----------------------------------------------------------------------
\section{Introduction}
%%----------------------------------------------------------------------

Until biffd owned the session record, a Claude session ended in its own process.
The process wrote a sentinel, tried to delete its row and release its tty name,
and a reaper in some other process finished what the first left.
\texttt{sentinel-reap.tex} models that mechanism and finds that a row can be left
with no process and no sentinel behind it ($BadOrphan$): a crash in one step (G1),
and the fallback sentinel racing its own reaper in five (G2).

With the record in the daemon the question changes.  The record is meant to be
durable (the code, today, keeps it in memory: $DaemonRestartForget$, below) and
the daemon, not a process that is going away, performs the end.  Two transitions
end a Claude session:

\begin{itemize}
  \item $EndOwnSession$: a connection of the \emph{current} run of the session
    ends its own session.  The daemon releases the tty name, deletes the row,
    appends a logout event and files the mail it holds for the session back into
    the user's inbox.  A connection of a run that a newer run has replaced is
    refused and changes nothing.
  \item $ExpireClaude$: a session with no connection that has been dormant for the
    retention period is removed with the same four effects.
\end{itemize}

Each of the four effects is its own step, and the record stays until all four are
done.  The obligations are:

\begin{enumerate}
  \item \textbf{No row outlives its record.}  The relay row and the tty name are
    held only while the record exists or is being ended, so there is no
    $BadOrphan$: the daemon has no state in which a row has nothing behind it.
  \item \textbf{No row outlives an end plus retention.}  A session nobody is
    attached to is eventually removed, once its retention is spent.
  \item \textbf{No mail is lost by an end.}  Mail issued is held by the session,
    still in the user's inbox, or shown; an end moves what it holds back to the
    inbox.
  \item \textbf{A clear ends the old session before the new hello.}  The new
    session's tty reservation is never released by the old session's end.
  \item \textbf{A replaced run cannot end the session.}
\end{enumerate}

%%----------------------------------------------------------------------
\section{Basic Types}
%%----------------------------------------------------------------------

\begin{zed}
  ZBOOL ::= ztrue | zfalse
\end{zed}

\begin{zed}
  RecState ::= rsNone | rsLive | rsDormant | rsEnding
\end{zed}

\begin{zed}
  TtyOwner ::= tNone | tOld | tNew
\end{zed}

\begin{axdef}
  maxAge : \nat \\
  maxMail : \nat
  \where
  maxAge = 2 \\
  maxMail = 2
\end{axdef}

$rsNone$ is no record (never started, or ended), $rsLive$ a record with a
connection, $rsDormant$ a record without one, and $rsEnding$ a record whose end
has begun and not finished.  $maxAge$ is the retention, in sweep ticks.
$tOld$ and $tNew$ are the old and the new session of a clear; both use the same
tty name, which is the reason the order matters.  Runs are $0$ and $1$.

%%----------------------------------------------------------------------
\section{State}
%%----------------------------------------------------------------------

\begin{schema}{Record}
  recState : RecState \\
  currentRun : 0 \upto 1 \\
  conns : \power (0 \upto 1) \\
  age : 0 \upto maxAge \\
  endRun : 0 \upto 1
\end{schema}

The daemon's record: its state, the current run of the session, the runs that
have a connection open (the connection of a replaced run is still open, and
refused), and the sweep ticks since the last connection left, and (a ghost) the run whose
connection began the end in progress.

\begin{schema}{Effects}
  rowPresent : ZBOOL \\
  ttyOwner : TtyOwner \\
  logged : ZBOOL \\
  rowTodo : ZBOOL \\
  ttyTodo : ZBOOL \\
  logTodo : ZBOOL \\
  mailTodo : ZBOOL \\
  newLive : ZBOOL \\
  ownEnded : ZBOOL
\end{schema}

What an end touches: the relay row, the tty-name reservation, the logout event,
and the four obligations an end in progress still owes.  $newLive$ is the new
session of a clear, after its hello.

\begin{schema}{Mail}
  outboxN : 0 \upto maxMail \\
  sessN : 0 \upto maxMail \\
  userN : 0 \upto maxMail \\
  shownN : 0 \upto maxMail \\
  issuedN : 0 \upto maxMail \\
  lostN : 0 \upto maxMail \\
  strandedN : 0 \upto maxMail \\
  shownDormant : ZBOOL \\
  forgot : ZBOOL
\end{schema}

Mail counts: held by the session ($outboxN$), in the session's own inbox in the
relay ($sessN$), in the user's mailbox ($userN$, where an end files what it
drains),
shown to the client ($shownN$), lost with a daemon that restarted ($lostN$, a
ghost), and issued in all ($issuedN$, a ghost).  $strandedN$ counts the mail left in the session inbox of a session whose
record a restart forgot ($sessN$ at that moment): nothing drains or reads it.
Mail already in the user's mailbox ($userN$) is not touched by the loss of a
record.  $shownDormant$ is a ghost
that records a print made while nobody was connected.  $forgot$ is a ghost set by
$DaemonRestartForget$, the restart that loses the record.

\begin{schema}{EndState}
  Record \\
  Effects \\
  Mail
\end{schema}

\begin{schema}{EndInv}
  EndState
  \where
  rowPresent = ztrue \implies ( recState \neq rsNone \lor forgot = ztrue ) \\
  ttyOwner = tOld \implies ( recState \neq rsNone \lor forgot = ztrue ) \\
  newLive = ztrue \implies ttyOwner = tNew \\
  recState = rsDormant \implies conns = \emptyset \\
  recState = rsLive \implies conns \neq \emptyset \\
  recState = rsEnding \implies endRun = currentRun \\
  recState \neq rsEnding \implies
    ( rowTodo = zfalse \land ttyTodo = zfalse \land logTodo = zfalse \land
      mailTodo = zfalse ) \\
  outboxN + sessN + userN + shownN + lostN + strandedN = issuedN \\
  shownDormant = zfalse
\end{schema}

The invariants, in the order of the obligations.  A row, or the old session's tty
name, exists only while the record does ($BadOrphan$ is not a state) unless the
daemon restarted and forgot the record ($forgot$, below).  The new
session of a clear always holds its tty name.  A live record has a connection and
a dormant one has none.  An end that is not in progress owes nothing.  And mail
is conserved: issued is held, in the session's inbox, in the user's mailbox,
shown, lost with a restart of the daemon (the outbox is its memory), or stranded
by a restart that forgot the record.  An end loses nothing: $FileOutbox$ files
what is held and drains the session's inbox into the user's mailbox.

\begin{schema}{Invariant}
  EndInv
\end{schema}

\begin{schema}{Init}
  EndState'
  \where
  recState' = rsLive \\
  currentRun' = 0 \\
  conns' = \{ 0 \} \\
  age' = 0 \\
  endRun' = 0 \\
  rowPresent' = ztrue \\
  ttyOwner' = tOld \\
  logged' = zfalse \\
  rowTodo' = zfalse \\
  ttyTodo' = zfalse \\
  logTodo' = zfalse \\
  mailTodo' = zfalse \\
  newLive' = zfalse \\
  ownEnded' = zfalse \\
  outboxN' = 0 \\
  sessN' = 0 \\
  userN' = 0 \\
  shownN' = 0 \\
  issuedN' = 0 \\
  lostN' = 0 \\
  strandedN' = 0 \\
  shownDormant' = zfalse \\
  forgot' = zfalse
\end{schema}

%%----------------------------------------------------------------------
\section{Operations on the record}
%%----------------------------------------------------------------------

\begin{schema}{ReplaceRun}
  \Delta Record \\
  \Xi Effects \\
  \Xi Mail
  \where
  recState = rsLive \\
  currentRun = 0 \\
  recState' = rsLive \\
  currentRun' = 1 \\
  conns' = conns \cup \{ 1 \} \\
  age' = age \\
  endRun' = endRun
\end{schema}

A resume with a fresh agent opens run $1$.  The connection of run $0$ stays open;
it belongs to a replaced run.

\begin{schema}{Detach}
  \Delta Record \\
  \Xi Effects \\
  \Xi Mail \\
  run? : 0 \upto 1
  \where
  recState = rsLive \\
  run? \in conns \\
  conns' = conns \setminus \{ run? \} \\
  recState' = ( \IF conns' = \emptyset \THEN rsDormant \ELSE rsLive ) \\
  currentRun' = currentRun \\
  age' = 0 \\
  endRun' = endRun
\end{schema}

A connection goes.  Detaching ends nothing: the record becomes dormant when its
last connection goes.

\begin{schema}{Reattach}
  \Delta Record \\
  \Xi Effects \\
  \Xi Mail
  \where
  recState = rsDormant \\
  recState' = rsLive \\
  currentRun' = currentRun \\
  conns' = \{ currentRun \} \\
  age' = 0 \\
  endRun' = endRun
\end{schema}

\begin{schema}{SweepTick}
  \Delta Record \\
  \Xi Effects \\
  \Xi Mail
  \where
  recState = rsDormant \\
  age < maxAge \\
  recState' = recState \\
  currentRun' = currentRun \\
  conns' = conns \\
  age' = age + 1 \\
  endRun' = endRun
\end{schema}

The daemon's sweep counts a dormant record's age.

\begin{schema}{DaemonRestartForget}
  \Delta Record \\
  \Xi Effects \\
  \Delta Mail
  \where
  recState \in \{ rsLive, rsDormant \} \\
  recState' = rsNone \\
  currentRun' = currentRun \\
  conns' = \emptyset \\
  age' = 0 \\
  endRun' = endRun \\
  outboxN' = 0 \\
  lostN' = lostN + outboxN \\
  sessN' = 0 \\
  userN' = userN \\
  strandedN' = strandedN + sessN \\
  shownN' = shownN \\
  issuedN' = issuedN \\
  shownDormant' = shownDormant \\
  forgot' = ztrue
\end{schema}

What the code does on a restart: biffd keeps the records of harness sessions in
memory, so the restart forgets them.  The relay row and the tty name stay, and age
out by the relay's own row lifetime ($RowAgesOut$, $TtyAgesOut$ below); the held
batch is lost with the outbox ($lostN$, PL-1753); and the mail waiting in the
session's own relay inbox ($sessN$) is neither drained into the user's mailbox
(that is part of an end) nor read by anyone, which the model counts as
$strandedN$.  It stays in the relay until the relay's own retention, and a
session resumed under the same id finds it again; the model does not state the
resume.  Mail already in the user's mailbox ($userN$) is untouched, and so is the
mail an end filed there.

This is not wholly safe, which is why the durable restart below is kept beside it
and not replaced by it: a forgotten session holds its tty name for the row
lifetime, and its unread mail is stranded.  The properties that mention the
record hold unless $forgot = ztrue$.  Persisting the records, or a scan of the
relay's rows when the daemon starts, would close both.

\begin{schema}{RowAgesOut}
  \Xi Record \\
  \Delta Effects \\
  \Xi Mail
  \where
  forgot = ztrue \\
  recState = rsNone \\
  rowPresent = ztrue \\
  rowPresent' = zfalse \\
  ttyOwner' = ttyOwner \\
  logged' = logged \\
  rowTodo' = rowTodo \\
  ttyTodo' = ttyTodo \\
  logTodo' = logTodo \\
  mailTodo' = mailTodo \\
  newLive' = newLive \\
  ownEnded' = ownEnded
\end{schema}

\begin{schema}{TtyAgesOut}
  \Xi Record \\
  \Delta Effects \\
  \Xi Mail
  \where
  forgot = ztrue \\
  recState = rsNone \\
  ttyOwner = tOld \\
  ttyOwner' = tNone \\
  rowPresent' = rowPresent \\
  logged' = logged \\
  rowTodo' = rowTodo \\
  ttyTodo' = ttyTodo \\
  logTodo' = logTodo \\
  mailTodo' = mailTodo \\
  newLive' = newLive \\
  ownEnded' = ownEnded
\end{schema}

The relay's row lifetime and the name reservation's lifetime run out.  This is the
NATS relay.  The local relay (no relay URL configured) keeps a reservation as a
lockfile with no lifetime: there $TtyAgesOut$ does not happen, and a forgotten
session holds its name until something releases it.  The model has no local
variant: the NATS behaviour is the favourable one.

\begin{schema}{DaemonRestart}
  \Delta Record \\
  \Xi Effects \\
  \Delta Mail
  \where
  recState \in \{ rsLive, rsDormant \} \\
  recState' = rsDormant \\
  currentRun' = currentRun \\
  conns' = \emptyset \\
  age' = ( \IF recState = rsLive \THEN 0 \ELSE age ) \\
  endRun' = endRun \\
  outboxN' = 0 \\
  lostN' = lostN + outboxN \\
  shownDormant' = shownDormant \\
  strandedN' = strandedN \\
  forgot' = forgot \\
  sessN' = sessN \\
  userN' = userN \\
  shownN' = shownN \\
  issuedN' = issuedN
\end{schema}

A restart, of a live record or of a dormant one (a detach leaves the outbox held
with nobody connected, and the daemon may die then): the record is durable and
returns dormant, with its age kept if it was already dormant; every connection
is gone.  The outbox is the daemon's memory, so what it held is lost from either
state (the window the unconfirmed batch has in DES-099, PL-1753); the mail the
relay holds is not.  A restart in the middle of an end ($rsEnding$) is not
modelled.

%%----------------------------------------------------------------------
\section{Ending a session}
%%----------------------------------------------------------------------

\begin{schema}{EndOwnSession}
  \Delta Record \\
  \Delta Effects \\
  \Xi Mail \\
  run? : 0 \upto 1
  \where
  recState = rsLive \\
  run? \in conns \\
  run? = currentRun \\
  recState' = rsEnding \\
  currentRun' = currentRun \\
  conns' = conns \\
  age' = age \\
  endRun' = run? \\
  rowPresent' = rowPresent \\
  ttyOwner' = ttyOwner \\
  logged' = logged \\
  rowTodo' = ztrue \\
  ttyTodo' = ztrue \\
  logTodo' = ztrue \\
  mailTodo' = ztrue \\
  newLive' = newLive \\
  ownEnded' = ztrue
\end{schema}

A connection of the current run ends its own session: the end begins and owes
its four effects.

\begin{schema}{EndRefused}
  \Xi EndState \\
  run? : 0 \upto 1
  \where
  recState = rsLive \\
  run? \in conns \\
  run? \neq currentRun
\end{schema}

A connection of a replaced run is refused and nothing changes.

\begin{schema}{ExpireClaude}
  \Delta Record \\
  \Delta Effects \\
  \Xi Mail
  \where
  recState = rsDormant \\
  conns = \emptyset \\
  age = maxAge \\
  recState' = rsEnding \\
  currentRun' = currentRun \\
  conns' = conns \\
  age' = age \\
  endRun' = currentRun \\
  rowPresent' = rowPresent \\
  ttyOwner' = ttyOwner \\
  logged' = logged \\
  rowTodo' = ztrue \\
  ttyTodo' = ztrue \\
  logTodo' = ztrue \\
  mailTodo' = ztrue \\
  newLive' = newLive \\
  ownEnded' = ownEnded
\end{schema}

Retention: a dormant record with no connection, aged to the retention, begins the
same end.

\begin{schema}{DeleteRow}
  \Xi Record \\
  \Delta Effects \\
  \Xi Mail
  \where
  recState = rsEnding \\
  rowTodo = ztrue \\
  rowPresent' = zfalse \\
  rowTodo' = zfalse \\
  ttyOwner' = ttyOwner \\
  logged' = logged \\
  ttyTodo' = ttyTodo \\
  logTodo' = logTodo \\
  mailTodo' = mailTodo \\
  newLive' = newLive \\
  ownEnded' = ownEnded
\end{schema}

\begin{schema}{ReleaseTty}
  \Xi Record \\
  \Delta Effects \\
  \Xi Mail
  \where
  recState = rsEnding \\
  ttyTodo = ztrue \\
  ttyOwner' = tNone \\
  ttyTodo' = zfalse \\
  rowPresent' = rowPresent \\
  logged' = logged \\
  rowTodo' = rowTodo \\
  logTodo' = logTodo \\
  mailTodo' = mailTodo \\
  newLive' = newLive \\
  ownEnded' = ownEnded
\end{schema}

The release is owner-checked, as the code's is: it frees the name only if this
session still holds it, and a name another session took stays theirs.  The
frozen variant \texttt{claude-session-end-unordered.tex} releases by name and
frees the new session's name; \texttt{claude-session-end-unordered-owned.tex}
keeps the owner check and lets the new hello run during the end, and the
invariant holds.  So the check, not the order alone, is what protects the new
session, and the order (the client waits for the end's reply) keeps the name
from being claimed while it is still held.

\begin{schema}{AppendLogout}
  \Xi Record \\
  \Delta Effects \\
  \Xi Mail
  \where
  recState = rsEnding \\
  logTodo = ztrue \\
  logged' = ztrue \\
  logTodo' = zfalse \\
  rowPresent' = rowPresent \\
  ttyOwner' = ttyOwner \\
  rowTodo' = rowTodo \\
  ttyTodo' = ttyTodo \\
  mailTodo' = mailTodo \\
  newLive' = newLive \\
  ownEnded' = ownEnded
\end{schema}

\begin{schema}{FileOutbox}
  \Xi Record \\
  \Delta Effects \\
  \Delta Mail
  \where
  recState = rsEnding \\
  mailTodo = ztrue \\
  userN' = userN + outboxN + sessN \\
  sessN' = 0 \\
  outboxN' = 0 \\
  shownN' = shownN \\
  issuedN' = issuedN \\
  lostN' = lostN \\
  shownDormant' = shownDormant \\
  strandedN' = strandedN \\
  forgot' = forgot \\
  mailTodo' = zfalse \\
  rowPresent' = rowPresent \\
  ttyOwner' = ttyOwner \\
  logged' = logged \\
  rowTodo' = rowTodo \\
  ttyTodo' = ttyTodo \\
  logTodo' = logTodo \\
  newLive' = newLive \\
  ownEnded' = ownEnded
\end{schema}

The mail the session holds is filed back in the relay and the session's inbox is
drained into the user's mailbox, where any reader finds it, and the end no longer
owes it.  A mail that arrives in the session's inbox after the drain and before
the record goes stays there, as in the code (the drain stops when a fetch finds
nothing and cannot see what comes later); the model allows the arrival and does
not state a property about it.

\begin{schema}{FinishEnd}
  \Delta Record \\
  \Xi Effects \\
  \Xi Mail
  \where
  recState = rsEnding \\
  rowTodo = zfalse \\
  ttyTodo = zfalse \\
  logTodo = zfalse \\
  mailTodo = zfalse \\
  recState' = rsNone \\
  currentRun' = currentRun \\
  conns' = \emptyset \\
  age' = age \\
  endRun' = endRun
\end{schema}

The record goes only when every effect is done.

%%----------------------------------------------------------------------
\section{A clear}
%%----------------------------------------------------------------------

\begin{schema}{NewHello}
  \Xi Record \\
  \Delta Effects \\
  \Xi Mail
  \where
  recState = rsEnding \\
  newLive = zfalse \\
  newLive' = ztrue \\
  ttyOwner' = tNew \\
  rowPresent' = rowPresent \\
  logged' = logged \\
  rowTodo' = rowTodo \\
  ttyTodo' = ttyTodo \\
  logTodo' = logTodo \\
  mailTodo' = mailTodo \\
  ownEnded' = ownEnded
\end{schema}

The new session's hello after a clear.  It is enabled only when the old session's
end is complete (no record, so no end in progress) and its tty name is free, and it
takes the name.  This is the ordering the client keeps by waiting for the end's
reply before it sends the hello.

%%----------------------------------------------------------------------
\section{Mail}
%%----------------------------------------------------------------------

\begin{schema}{MailArrive}
  \Xi Record \\
  \Xi Effects \\
  \Delta Mail
  \where
  issuedN < maxMail \\
  sessN' = sessN + 1 \\
  userN' = userN \\
  issuedN' = issuedN + 1 \\
  outboxN' = outboxN \\
  shownN' = shownN \\
  lostN' = lostN \\
  shownDormant' = shownDormant \\
  strandedN' = strandedN \\
  forgot' = forgot
\end{schema}

\begin{schema}{MailFetch}
  \Xi Record \\
  \Xi Effects \\
  \Delta Mail
  \where
  recState = rsLive \\
  sessN > 0 \\
  sessN' = sessN - 1 \\
  userN' = userN \\
  outboxN' = outboxN + 1 \\
  shownN' = shownN \\
  issuedN' = issuedN \\
  lostN' = lostN \\
  shownDormant' = shownDormant \\
  strandedN' = strandedN \\
  forgot' = forgot
\end{schema}

The client prints a message it was handed.  It needs a connected client: with the
session dormant (a daemon restart left no connection) the held mail cannot be
shown, and a reply already in flight when the daemon died counts as shown, which
is the count $shownN$ took when it was handed over.

\begin{schema}{MailShow}
  \Xi Record \\
  \Xi Effects \\
  \Delta Mail
  \where
  recState = rsLive \\
  outboxN > 0 \\
  outboxN' = outboxN - 1 \\
  shownN' = shownN + 1 \\
  sessN' = sessN \\
  userN' = userN \\
  issuedN' = issuedN \\
  lostN' = lostN \\
  shownDormant' = ( \IF recState = rsDormant \THEN ztrue \ELSE shownDormant ) \\
  strandedN' = strandedN \\
  forgot' = forgot
\end{schema}

\begin{schema}{Quiesce}
  \Xi EndState
\end{schema}

%%----------------------------------------------------------------------
\section{Model Checking}
%%----------------------------------------------------------------------

$EndInv$ is the state invariant.  The retention obligation is stated as
\emph{enabledness}, plus one reachability check, and not as liveness.  $EF$
alone would be satisfied by a client that reattaches and ends the session, which
is not retention, and $AF$ cannot hold of a machine whose $Quiesce$ may be taken
forever.  So:

\begin{quote}
  $AG\,(\,recState = rsDormant \land age = maxAge \Rightarrow
  enabled(ExpireClaude)\,)$ \\
  $AG\,(\,recState = rsEnding \land rowTodo = ztrue \Rightarrow
  enabled(DeleteRow)\,)$, and the same for $ReleaseTty$, $AppendLogout$ and
  $FileOutbox$ with their flags \\
  $AG\,(\,recState = rsEnding \land rowTodo = ttyTodo = logTodo = mailTodo =
  zfalse \Rightarrow enabled(FinishEnd)\,)$ \\
  $AG\,(\,recState = rsEnding \Rightarrow EF\;recState = rsNone\,)$, a
  \emph{reachability} check: from an end in progress the record can be removed
\end{quote}

Under the assumption that the daemon's sweep and the end's steps are scheduled
fairly (a step that stays enabled is eventually taken, and $Quiesce$ stands for the
absence of work, so it is taken only when nothing else is enabled), the
enabledness formulas give: a dormant session at its retention is removed, and an
end in progress finishes.  The model does not state fairness, and none of the
formulas is a liveness theorem.

and the non-vacuity goals are that both transitions are reachable:
$recState = rsEnding$ after $EndOwnSession$ and after $ExpireClaude$, that
$EndRefused$ is enabled ($currentRun = 1 \land 0 \in conns$), and that a clear
completes through an end ($newLive = ztrue \land ownEnded = ztrue \land forgot =
zfalse$, the ghost $ownEnded$ set by $EndOwnSession$: a trace with an
$EndOwnSession$ step before $NewHello$, not the path through the forgetting restart
and $TtyAgesOut$, which enables $NewHello$ with no end).

Eight frozen variants are the negative controls, each the shipped file with one
thing taken out:
\begin{itemize}
  \item \texttt{claude-session-end-nofile.tex}: $FileOutbox$ drops what the session
    holds (the end that does not file mail back, the requirement A2 of DES-099
    states): the conservation invariant is violated.
  \item \texttt{claude-session-end-unordered.tex}: $NewHello$ is enabled while the
    old end is still in progress and $ReleaseTty$ frees the name by name: the old
    end frees the new session's name, and $newLive = ztrue \Rightarrow ttyOwner =
    tNew$ is violated.
  \item \texttt{claude-session-end-unordered-owned.tex}: the same unordered
    $NewHello$ with the owner-checked $ReleaseTty$ of the shipped model: the
    invariant \textbf{holds} over its 2,351 states.  This is the positive control:
    the owner check is what keeps the new session's name.
  \item \texttt{claude-session-end-orphan.tex}: $ExpireClaude$ owes no row deletion
    (the retention path forgets the row): the record goes with the row still
    present, and $rowPresent = ztrue \Rightarrow recState \neq rsNone$ is violated:
    $BadOrphan$.
  \item \texttt{claude-session-end-noexpire.tex}: the guard of $ExpireClaude$ is
    $age > maxAge$, which the saturating sweep never reaches: the first retention
    formula is \textbf{FALSE}, a dormant record at its retention is never
    removed.
  \item \texttt{claude-session-end-showanywhere.tex}: $MailShow$ needs no
    connection.  Held mail is consumed by a print with nobody connected:
    $shownDormant = zfalse$ is violated, by $MailArrive$, $MailFetch$, $Detach$,
    $MailShow$.  (A restart is not the route: it loses the outbox in this control
    as in the shipped model, so the held mail has nothing left to show.)
  \item \texttt{claude-session-end-anyrun.tex}: $EndOwnSession$ accepts a
    connection of a replaced run: the replaced connection ends the session the
    newer run owns (checked by the goal $recState = rsEnding$ reached with
    $run? \neq currentRun$, as the history of the operation).
\end{itemize}

%%----------------------------------------------------------------------
%% Reproducing the checks.  Run from the repo root; one probcli at a time.
%%
%%   PROBCLI="${PROBCLI:-$HOME/Applications/ProB/probcli}"
%%   P="-p DEFAULT_SETSIZE 1 -p MAXINT 6 -p MAX_INITIALISATIONS 20 \
%%      -p MAX_OPERATIONS 500 -p TIME_OUT 600000"
%%   ( ulimit -v 8000000; "$PROBCLI" docs/specs/claude-session-end.tex $P -model_check )
%%   "$PROBCLI" docs/specs/claude-session-end.tex $P -model_check -ctlformula \
%%       "AG({recState=rsDormant & age=2} => e(ExpireClaude))"
%%   "$PROBCLI" docs/specs/claude-session-end.tex $P -model_check -ctlformula \
%%       "AG({recState=rsEnding & rowTodo=ztrue} => e(DeleteRow))"
%%   "$PROBCLI" docs/specs/claude-session-end.tex $P -model_check -ctlformula \
%%       "AG({recState=rsEnding & rowTodo=zfalse & ttyTodo=zfalse & logTodo=zfalse & mailTodo=zfalse} => e(FinishEnd))"
%%   "$PROBCLI" docs/specs/claude-session-end.tex $P -model_check -ctlformula \
%%       "AG({recState=rsEnding} => EF {recState=rsNone})"   (reachability)
%%   goals: "recState=rsEnding & age=2", "newLive=ztrue & ownEnded=ztrue & forgot=zfalse", "recState=rsLive & conns/={} & currentRun=1"
%%   "$PROBCLI" docs/specs/claude-session-end.tex $P -model_check -ctlformula \
%%       "AG({forgot=ztrue & recState=rsNone & rowPresent=ztrue} => e(RowAgesOut))"   (and TtyAgesOut)
%%   goals (the restart that forgets): "forgot=ztrue & rowPresent=ztrue", "strandedN>0",
%%       "forgot=ztrue & ttyOwner=tNew & rowPresent=zfalse"
%%   shipped model, exhaustive: 1,411 states; unordered-owned: 2,351 states, no violation
%%
%% Negative controls (each reports an invariant violation, the noexpire one a
%% FALSE CTL formula):
%%
%%   for V in nofile unordered orphan anyrun showanywhere; do
%%     "$PROBCLI" docs/specs/claude-session-end-$V.tex $P -model_check
%%   done
%%   "$PROBCLI" docs/specs/claude-session-end-noexpire.tex $P -model_check -ctlformula \
%%       "AG({recState=rsDormant & age=2} => e(ExpireClaude))"
%%----------------------------------------------------------------------

\end{document}
